\documentclass[fontsize=11pt, parskip=half]{scrartcl}
\usepackage[paperwidth=6.5in, paperheight=9.5in,
            margin=1cm]{geometry}
\pagestyle{plain}

\usepackage{amsmath, amssymb, amsthm, amsfonts}

\usepackage[dvipsnames]{xcolor}

\usepackage[draft]{commenting}
\declareauthor{leon}{LL}{blue}
\authorcommand{leon}{comment}

\usepackage[colorlinks=true, linkcolor=blue, urlcolor=blue, citecolor=blue]{hyperref}
\usepackage{url}
\usepackage[capitalise,noabbrev]{cleveref}

% 3. Define your theorem environments AFTER cleveref is loaded
\newtheorem{theorem}{Theorem}[section]
\newtheorem{lemma}[theorem]{Lemma}
\newtheorem{fact}[theorem]{Fact}
\theoremstyle{definition}
\newtheorem{definition}[theorem]{Definition}
\newtheorem{remark}[theorem]{Remark}

% 4. Set up cross-reference names
\crefname{section}{Problem}{Problems}
\Crefname{section}{Problem}{Problems}
\crefname{subsection}{Problem}{Problems}
\Crefname{subsection}{Problem}{Problems}

\crefname{theorem}{Theorem}{Theorems}
\Crefname{theorem}{Theorem}{Theorems}

\crefname{definition}{Definition}{Definitions}
\Crefname{definition}{Definition}{Definitions}

\crefname{lemma}{Lemma}{Lemmas}
\Crefname{lemma}{Lemma}{Lemmas}

\crefname{fact}{Fact}{Facts}
\Crefname{fact}{Fact}{Facts}

% Solution toggle: set to true to show solutions, false to hide
\newif\ifsolutions
\solutionstrue  % comment out or set \solutionsfalse to hide

\usepackage{tcolorbox}
\tcbuselibrary{breakable}

\ifsolutions
  \newenvironment{solution}{%
    \begin{tcolorbox}[colback=black!5!white, colframe=black!50!white,
      title={\textbf{Solution}}, fonttitle=\normalfont, breakable]
  }{\end{tcolorbox}}
\else
  \newenvironment{solution}{\expandafter\comment}{\expandafter\endcomment}
  \usepackage{verbatim} % for \comment environment
\fi

\tcolorboxenvironment{definition}{colback=blue!3!white, colframe=blue!40!white,
  breakable, top=4pt, bottom=4pt}
\tcolorboxenvironment{theorem}{colback=red!3!white, colframe=red!40!white,
  breakable, top=4pt, bottom=4pt}
\tcolorboxenvironment{lemma}{colback=purple!3!white, colframe=purple!40!white,
  breakable, top=4pt, bottom=4pt}
\tcolorboxenvironment{fact}{colback=orange!3!white, colframe=orange!40!white,
  breakable, top=4pt, bottom=4pt}

\newtcolorbox{defbox}[1][]{colback=blue!3!white, colframe=blue!40!white,
  title={\textbf{#1}}, fonttitle=\normalfont, breakable, top=4pt, bottom=4pt}

\newtcolorbox{examplebox}[1][]{colback=green!3!white, colframe=green!40!black,
  title={\textbf{#1}}, fonttitle=\normalfont, breakable, top=4pt, bottom=4pt}
%%%%%%%%%%%%%%%%%%%%%%%%%

\newcommand{\mytitle}[1]{{\noindent\Large\textbf{#1}}}

\renewcommand{\thesection}{\arabic{section}}

\newcommand{\mysubsection}[2]{\par\medskip\noindent\refstepcounter{subsection}\textbf{\thesubsection.}~#2\par\smallskip}

\newcommand{\Msol}{\mathcal{M}_{sol}}
\newcommand{\Mc}{\mathcal{M}}
\newcommand{\KL}{D_{\text{KL}}}
\newcommand{\Dn}[2]{D_{n}(#1 \,\|\, #2)}
\newcommand{\dt}[2]{d_{t}(#1 \,\|\, #2)}
\newcommand{\Sn}[2]{S_{n}(#1 \,\|\, #2)}

%===================================
\begin{document}

\noindent\textbf{Iliad Intensive} \hfill \textbf{London Initiative for Safe AI}\\
April 2026 \hfill Leon Lang\\

\mytitle{Solomonoff Induction Exercises \hfill \today}

\medskip
\noindent
Much of this material is drawn from Hutter's \emph{An Introduction to Universal Artificial Intelligence}~\cite{Hutter:24uaibook2} (Chapter~3 in particular) and the earlier \emph{Universal Artificial Intelligence}~\cite{Hutter:04uaibook}; both books provide fuller explanations, additional context, and proofs that we gloss over here.

\medskip
\noindent
\textbf{Difficulty ratings.} Each problem is tagged with a rating in square brackets following Knuth's exercise rating scheme~\cite{Knuth:73a}: roughly, \textbf{[00]}~trivial, \textbf{[10]}~15--minute pencil-and-paper, \textbf{[20]}~1--2~hours, \textbf{[30]}~several hours to a day, \textbf{[40]}~a significant research result. Intermediate values are possible.

%===================================
\section*{Goal and Roadmap}

\textbf{A recipe for prediction:} Solomonoff induction is an attempt to mathematically
formalize the \href{https://en.wikipedia.org/wiki/Problem_of_induction}\emph{problem of induction}
from philosophy: How to make predictions about the future based on past observations?

We formalize this as sequence prediction: There is some true environment $\mu$ that generates
a sequence $x_1, x_2, \dots$ of (binary) symbols. Each next symbol $x_t \sim \mu(\cdot \mid x_{<t})$ is sampled
from $\mu$ conditioned on the past $x_{<t}$. A predictor $P$ takes the history $x_{<t}$ and gives a 
distribution $P(\cdot \mid x_{<t})$ over the next symbol $x_t$. 

We measure the quality of a predictor $P$ by the $\mu$-expected squared prediction error, summed over every timestep:
\begin{align}
S_\infty^\mu := \sum_{t=1}^{\infty} \sum_{x_{<t} \in \mathbb{B}^*} \mu(x_{<t}) \sum_{x_t \in \mathbb{B}} \big( P(x_t \mid x_{<t}) - \mu(x_t \mid x_{<t}) \big)^2.
\end{align}

How to construct a predictor $P$ such that $S^\mu_\infty$ is small (or at least finite)? Bayesian inference to the rescue: we choose as our predictor a \emph{Bayesian mixture} $\xi$ over a countable class $\Mc = \{\nu_1, \nu_2, \dots\}$ of candidate environments, weighted by prior beliefs $w_\nu$ (formalized in \cref{def:mixture}). The main results we build up to are:
\begin{itemize}\itemsep=2pt
  \item \textbf{Cumulative bound} (\cref{prob:bound}): Assuming $\mu \in \Mc$, $S^\mu_\infty \leq -\ln w_\mu$. The higher the prior
  on $\mu$, the lower the prediction error.
  \item \textbf{Explicit bound} (\cref{prob:bound_explicit}): specialize to the Solomonoff prior $w_\nu = 2^{-K(\nu)}$ to get $S^\mu_\infty \leq K(\mu)\ln 2$, where $K$ is the \emph{Kolmogorov complexity}.
  \item \textbf{Pareto optimality} (\cref{prob:kl_pareto,prob:sq_pareto}): no other predictor weakly dominates $\xi$ on every $\nu \in \Mc$, for either KL or squared loss.
  \item \textbf{Misspecified version} (\cref{prob:misspec}): If $\mu \not\in \Mc$, $S^\mu_\infty$ diverges at a rate 
  $O(\KL(\hat{\mu}||\mu))$, where $\hat{\mu}$ is the ``closest'' environment to $\mu$ with $\hat{\mu} \in \Mc$.
\end{itemize}

%===================================
\section*{Notation}

For more background, see \href{https://www.lesswrong.com/posts/HSDumToH57nSRdLST/a-technical-introduction-to-solomonoff-induction-without-k}{the corresponding post on Solomonoff induction}.

\begin{defbox}[Symbols used throughout]
\begin{itemize}\itemsep=2pt
  \item $\mathbb{B} = \{\texttt{0}, \texttt{1}\}$ --- the binary alphabet; $\mathbb{B}^*$ all finite binary strings; $\mathbb{B}^n$ length-$n$ strings.
  \item $\epsilon$ --- the empty string. $x_{1:n} := x_1 x_2 \cdots x_n$; $x_{<t} := x_1 x_2 \cdots x_{t-1}$.
  \item $X_{1:n}, X_{<t}$ --- the corresponding random variables.
  \item $\nu, \rho$ --- generic environments / predictors. $\mu$ --- the true (unknown) environment generating the data. $\xi$ --- a Bayesian mixture of environments.
  \item $\Mc = \{\nu_1, \nu_2, \ldots\}$ --- a countable class of candidate environments.
  \item $w_\nu > 0$ with $\sum_{\nu \in \Mc} w_\nu = 1$ --- the prior over $\Mc$. $w(\nu \mid x_{<t})$ --- the posterior after observing $x_{<t}$.
  \item $\Delta \mathcal{X}$ --- the set of probability distributions over a finite set $\mathcal{X}$.
\end{itemize}
\end{defbox}

\begin{defbox}[The Solomonoff prior]
All results in this sheet hold for any countable class $\Mc$ and prior $w_\nu$. The canonical \textbf{Solomonoff} choice takes $\Mc$ to be the class of all computable environments and the \textbf{Solomonoff prior}
\[
w_\nu ~:=~ 2^{-K(\nu)},
\]
where $K(\nu)$ is the \emph{Kolmogorov complexity} of $\nu$: the length of the shortest program that computes $\nu$. The Bayesian mixture $\xi$ under this particular prior is the \emph{universal mixture}, written $\xi_U$. With this class the assumption $\mu \in \Mc$ reduces to ``the universe is computable'' --- as weak an assumption as one can make. We do not make use of
details of defining $K$ or computability, and dodge it for today.\footnote{\textit{Technical aside:} The class of all computable measures is countable but not enumerable (by the halting problem, you can't decide which Turing machines define total functions), so a faithful construction of $\xi_U$ requires expanding $\Mc$ to all \emph{lower-semicomputable semimeasures} (Levin), at which point $\xi_U$ becomes a semimeasure rather than a measure --- $\sum_x \xi_U(x) \leq 1$ instead of $=1$, conditionals may not sum to $1$, and every other proof in this sheet gains a side-case. We dodge all of this: $\Mc$ is just \emph{some} countable set of proper measures with $\mu \in \Mc$, and we assume $K(\nu)$ is defined for each $\nu \in \Mc$ when we need it (\cref{prob:bound_explicit,prob:dominance,prob:error_count,prob:mistakes}). We never need $\xi$ itself to be in $\Mc$, nor do we need its computability, so we don't ask. See \cite[\S2.4.3]{Hutter:04uaibook} for the full Levin construction.}
\end{defbox}

%===================================
\section{The mixture is a predictor}\label{prob:one_step}

\begin{definition}[Environment]\label{def:environment}
An \textbf{environment} $\nu$ assigns to each history $x_{<t} \in \mathbb{B}^*$ a predictive distribution $\nu(\cdot \mid x_{<t}) \in \Delta \mathbb{B}$ over the next symbol. The joint probability of a string is defined by the chain rule
\[
\nu(x_{1:n}) ~:=~ \prod_{t=1}^{n} \nu(x_t \mid x_{<t}), \qquad \nu(\epsilon) := 1,
\]
so that $\nu(x_{1:t}) = \nu(x_{<t}) \cdot \nu(x_t \mid x_{<t})$ and $\nu(x) = \nu(x\texttt{0}) + \nu(x\texttt{1})$ for all $x \in \mathbb{B}^*$. The true (unknown) environment generating the data is denoted $\mu$.
\end{definition}

\begin{definition}[Bayesian mixture $\xi$]\label{def:mixture}
Let $\Mc = \{\nu_1, \nu_2, \dots\}$ be a countable class of environments with prior weights $w_\nu > 0$ satisfying $\sum_{\nu \in \Mc} w_\nu = 1$. The \textbf{Bayesian mixture} is defined as a prior-weighted mixture over all environments $\nu$ in $\Mc$:
\begin{equation}\label{eq:mixture_expression}
  \xi(x_{1:t}) ~:=~ \sum_{\nu \in \Mc} w_\nu\, \nu(x_{1:t}).
\end{equation}
Its one-step predictive distribution is
\[
  \xi(x_t \mid x_{<t}) ~=~ \sum_{\nu \in \Mc} w(\nu \mid x_{<t})\, \nu(x_t \mid x_{<t}),
\]
with posterior weights 
\[
w(\nu \mid x_{<t}) ~:=~ w_\nu\, \frac{\nu(x_{<t})}{\xi(x_{<t})}
\qquad w(\nu \mid \epsilon) := w_\nu.
\]
Here $w(\nu \mid x_{<t})$ is the \emph{posterior} belief in $\nu$ after observing $x_{<t}$.
\end{definition}

The mixture $\xi$ is defined as a sum of \emph{joint} probabilities. To use it as a predictor we need its one-step conditional $\xi(x_t \mid x_{<t})$, and we need to know it is a genuine probability distribution (so that, later, Pinsker's inequality applies to it).

\mysubsection{1}{[10] (Mixture is a predictor)\label{prob:one_step_ex}}

Starting from $\xi(x_{1:t}) = \sum_{\nu \in \Mc} w_\nu\, \nu(x_{1:t})$, show that
\[
\xi(x_t \mid x_{<t}) ~=~ \sum_{\nu \in \Mc} w(\nu \mid x_{<t})\, \nu(x_t \mid x_{<t}),
\qquad w(\nu \mid x_{<t}) := w_\nu\, \frac{\nu(x_{<t})}{\xi(x_{<t})},
\]
and conclude that $\sum_{x_t \in \mathbb{B}} \xi(x_t \mid x_{<t}) = 1$, i.e.\ $\xi(\cdot \mid x_{<t})$ is a probability distribution over $\mathbb{B}$.

\textit{Hint:} Write $\xi(x_t \mid x_{<t}) = \xi(x_{1:t}) / \xi(x_{<t})$, expand the numerator, and use the chain rule $\nu(x_{1:t}) = \nu(x_{<t})\, \nu(x_t \mid x_{<t})$ (\cref{def:environment}). For the last part, note that the posterior weights sum to $1$.

\begin{solution}
The one-step conditional is the ratio of joint to marginal, $\xi(x_t \mid x_{<t}) := \xi(x_{1:t}) / \xi(x_{<t})$. Expanding the numerator and applying the chain rule $\nu(x_{1:t}) = \nu(x_{<t})\, \nu(x_t \mid x_{<t})$ to each term:
\begin{align*}
\xi(x_t \mid x_{<t})
~&=~ \frac{\sum_{\nu \in \Mc} w_\nu\, \nu(x_{1:t})}{\xi(x_{<t})}
~=~ \frac{\sum_{\nu \in \Mc} w_\nu\, \nu(x_{<t})\, \nu(x_t \mid x_{<t})}{\xi(x_{<t})} \\
~&=~ \sum_{\nu \in \Mc} \underbrace{\frac{w_\nu\, \nu(x_{<t})}{\xi(x_{<t})}}_{=\, w(\nu \mid x_{<t})}\, \nu(x_t \mid x_{<t})
~=~ \sum_{\nu \in \Mc} w(\nu \mid x_{<t})\, \nu(x_t \mid x_{<t}).
\end{align*}
The posterior weights are non-negative and sum to $1$:
\[
\sum_{\nu \in \Mc} w(\nu \mid x_{<t})
~=~ \frac{\sum_{\nu \in \Mc} w_\nu\, \nu(x_{<t})}{\xi(x_{<t})}
~=~ \frac{\xi(x_{<t})}{\xi(x_{<t})}
~=~ 1.
\]
Therefore $\xi(\cdot \mid x_{<t})$ is a convex combination of the probability distributions $\nu(\cdot \mid x_{<t})$, so it is itself a probability distribution over $\mathbb{B}$:
\[
\sum_{x_t \in \mathbb{B}} \xi(x_t \mid x_{<t})
~=~ \sum_{\nu \in \Mc} w(\nu \mid x_{<t}) \underbrace{\sum_{x_t \in \mathbb{B}} \nu(x_t \mid x_{<t})}_{=\, 1}
~=~ \sum_{\nu \in \Mc} w(\nu \mid x_{<t})
~=~ 1. \qedhere
\]
\end{solution}

%===================================
\section{Chain rule for KL divergence}\label{prob:chain_rule}

\begin{definition}[Kullback--Leibler divergence]\label{def:kl}
Let $P$ and $Q$ be probability distributions over a finite set $\mathcal{X}$. The \textbf{Kullback--Leibler (KL) divergence} from $P$ to $Q$ is
\[
D_{\text{KL}}(P \,\|\, Q) := \sum_{x \in \mathcal{X}} P(x) \ln \frac{P(x)}{Q(x)},
\]
where we define $0 \ln \frac{0}{q} := 0$ for any $q \geq 0$, and $p \ln \frac{p}{0} := \infty$ for $p > 0$.
\end{definition}

\begin{definition}[Joint and conditional KL divergence]\label{def:joint_kl}
Let $P$ and $Q$ be joint probability distributions over random variables $X, Y$ taking values in $\mathcal{X}, \mathcal{Y}$. The \textbf{joint} and \textbf{conditional} KL divergences are
\[
D_{\text{KL}}(P(X,Y) \,\|\, Q(X,Y)) := \sum_{x, y} P(x,y) \ln \frac{P(x,y)}{Q(x,y)},
\qquad
D_{\text{KL}}(P(Y \mid X) \,\|\, Q(Y \mid X)) := \sum_{x} P(x) \sum_{y} P(y \mid x) \ln \frac{P(y \mid x)}{Q(y \mid x)},
\]
with the same conventions as \cref{def:kl}.
\end{definition}

\mysubsection{1}{[10] (KL chain rule)\label{prob:chain_rule_ex}}

We show that KL divergence satisfies the chain rule:
\[
D_{\text{KL}}(P(X,Y) \,\|\, Q(X,Y))
=
D_{\text{KL}}(P(X) \,\|\, Q(X))
+
D_{\text{KL}}(P(Y \mid X) \,\|\, Q(Y \mid X)).
\]

\begin{solution}
We expand the joint KL divergence using the standard chain rule for probabilities, $P(x,y) = P(x) P(y \mid x)$:
\begin{align*}
& D_{\text{KL}}(P(X,Y) \,\|\, Q(X,Y)) \\
&= \sum_{x,y} P(x,y)
\ln \frac{P(x,y)}{Q(x,y)} \\
&= \sum_{x,y} P(x) P(y \mid x)
\ln \frac{P(x) P(y \mid x)}{Q(x) Q(y \mid x)} \\
&= \sum_{x} \sum_y P(x) P(y \mid x)
\left[
\ln \frac{P(x)}{Q(x)}
+
\ln \frac{P(y \mid x)}{Q(y \mid x)}
\right] \\
&= \sum_x P(x) \ln \frac{P(x)}{Q(x)}
\underbrace{\sum_y P(y \mid x)}_{=1}
 + \sum_x P(x) \sum_y P(y \mid x)
\ln \frac{P(y \mid x)}{Q(y \mid x)} \\
&= D_{\text{KL}}(P(X) \,\|\, Q(X))
+
D_{\text{KL}}(P(Y \mid X) \,\|\, Q(Y \mid X)).
\qedhere
\end{align*}
\end{solution}

%===================================
\section{Telescoping the chain rule}\label{prob:telescope}

The two-variable chain rule of \cref{prob:chain_rule} extends, by induction, to a sequence of variables. The point of this telescoping identity is that it lets us trade a single joint KL divergence over an entire history for a sum of per-step conditional KL divergences (and vice versa)---exactly the bridge we will need in \cref{prob:bound}.

We introduce two pieces of notation that will be used throughout the rest of the worksheet:

\begin{definition}[Cumulative and per-step KL]\label{def:Dn}
For environments / predictors $\nu, \rho$ and horizon $n \geq 1$, the \textbf{cumulative joint KL} and the \textbf{expected per-step KL at step $t$} are:
\begin{align*}
  \Dn{\nu}{\rho} ~&:=~ \KL\!\bigl(\nu(X_{1:n}) \,\big\|\, \rho(X_{1:n})\bigr) \\
                  &\phantom{:}=~ \sum_{x_{1:n} \in \mathbb{B}^n}
                       \nu(x_{1:n}) \,\ln \frac{\nu(x_{1:n})}{\rho(x_{1:n})}, \\[6pt]
  \dt{\nu}{\rho} ~&:=~ \mathbb{E}_{X_{<t} \sim \nu}\!\Bigl[\KL\!\bigl(\nu(\cdot \mid X_{<t}) \,\big\|\, \rho(\cdot \mid X_{<t})\bigr)\Bigr] \\
                  &\phantom{:}=~ \sum_{x_{<t} \in \mathbb{B}^{t-1}}
                       \nu(x_{<t}) \sum_{x_t \in \mathbb{B}}
                       \nu(x_t \mid x_{<t}) \,\ln \frac{\nu(x_t \mid x_{<t})}{\rho(x_t \mid x_{<t})}.
\end{align*}
The first argument always names the distribution the expectation is taken over (we do \emph{not} adopt Hutter's convention of an implicit ``$\mathbb{E}$ always over $\mu$,'' since the Pareto sections need $\nu$ to vary across the model class). \Cref{prob:telescope} below establishes that the two are related by $\Dn{\nu}{\rho} = \sum_{t=1}^n \dt{\nu}{\rho}$.
\end{definition}

\mysubsection{1}{[10] (Telescoping the chain rule)\label{prob:telescope_ex}}

Show by induction on $n$ (using \cref{prob:chain_rule}) that for all environments $\nu, \rho$,
\[
\Dn{\nu}{\rho} ~=~ \sum_{t=1}^{n} \dt{\nu}{\rho},
\]
where on the right-hand side, the $t=1$ summand uses the conventions $\nu(X_1 \mid X_{<1}) := \nu(X_1)$, $\rho(X_1 \mid X_{<1}) := \rho(X_1)$.

\begin{solution}
We prove this by induction on $n$. Unpacking the per-step term via \cref{def:joint_kl}, $\dt{\nu}{\rho} = D_{\text{KL}}(\nu(X_t \mid X_{<t}) \,\|\, \rho(X_t \mid X_{<t}))$, so the claim is equivalent to $D_{\text{KL}}(\nu(X_{1:n}) \,\|\, \rho(X_{1:n})) = \sum_{t=1}^n D_{\text{KL}}(\nu(X_t \mid X_{<t}) \,\|\, \rho(X_t \mid X_{<t}))$.

\textbf{Base case ($n = 1$):} The sum reduces to $D_{\text{KL}}(\nu(X_1 \mid X_{<1}) \,\|\, \rho(X_1 \mid X_{<1})) = D_{\text{KL}}(\nu(X_1) \,\|\, \rho(X_1))$, which is just the left-hand side.

\textbf{Induction step:} Apply the chain rule from \cref{prob:chain_rule} with $X = X_{<n}$ and $Y = X_n$:
\begin{align*}
D_{\text{KL}}(\nu(X_{1:n}) \,\|\, \rho(X_{1:n}))
&= D_{\text{KL}}(\nu(X_{<n}) \,\|\, \rho(X_{<n})) + D_{\text{KL}}(\nu(X_n \mid X_{<n}) \,\|\, \rho(X_n \mid X_{<n})).
\end{align*}
By the induction hypothesis, the first term satisfies
\[
D_{\text{KL}}(\nu(X_{<n}) \,\|\, \rho(X_{<n})) = \sum_{t=1}^{n-1} D_{\text{KL}}(\nu(X_t \mid X_{<t}) \,\|\, \rho(X_t \mid X_{<t})).
\]
Combining both gives the claim. \qedhere
\end{solution}

%===================================
\section{Pinsker's inequality}\label{prob:pinsker}

The next ingredient relates squared prediction error to KL divergence~\cite[Lemma~11.6.1]{Cover:06}. We state the form we will actually use and defer its proof to \cref{app:pinsker}. The exercise here is the key analytic step: the binary inequality.

\begin{theorem}[Pinsker's inequality]\label{thm:pinsker}
Let $P$ and $Q$ be probability distributions over $\mathbb{B} = \{\texttt{0}, \texttt{1}\}$, i.e., $q_0, q_1 \geq 0$
with $q_0 + q_1 = 1$. Write $p_x = P(x)$ and $q_x = Q(x)$. Then
\[
    \sum_{x \in \mathbb{B}} \big( q_x - p_x \big)^2 \;\leq\; \sum_{x \in \mathbb{B}} p_x \ln \frac{p_x}{q_x}.
    \tag{$\star$}
\]
The proof is deferred to \cref{app:pinsker}; it reduces to the binary inequality of \cref{lem:pinsker_binary} below (plus the boundary case $q_x = 0$).
\end{theorem}

\begin{lemma}[Pinsker's binary inequality]\label{lem:pinsker_binary}
For $0 \leq p \leq 1$ and $0 < q < 1$ we have
\[
    2(p-q)^2 ~\leq~ p \ln \frac{p}{q} + (1-p) \ln \frac{1-p}{1-q}.
\]
\end{lemma}

\mysubsection{1}{[15] (Pinsker binary inequality)\label{prob:pinsker_ex}}

Prove \cref{lem:pinsker_binary}.

\textit{Hint:} Define a function $f(x)$ such that the right-hand side can be written as $f(p) - f(q)$ (try seperating the LHS into terms with and without $q$). Then use the fundamental theorem of calculus, $f(p) - f(q) = \int_q^p f'(x) \, dx$, together with the inequality $x(1-x) \leq 1/4$ for $0 \leq x \leq 1$, to obtain the result.

\begin{solution}
Define $f(x) := p \ln x + (1-p) \ln (1-x)$, so that
\[
p \ln \frac{p}{q} + (1-p) \ln \frac{1-p}{1-q} = f(p) - f(q)
= \int_q^p f'(x)\, dx
= \int_q^p \frac{p-x}{x(1-x)}\, dx,
\]
using $f_p'(x) = \frac{p}{x} - \frac{1-p}{1-x} = \frac{p-x}{x(1-x)}$. Recall also that $x(1-x) \leq \tfrac14$ on $[0,1]$, so $\frac{1}{x(1-x)} \geq 4$.

\textbf{Case $q \leq p$.} For every $x \in [q,p]$ we have $x \leq p$, hence $p - x \geq 0$. Multiplying $\frac{1}{x(1-x)} \geq 4$ by the nonnegative factor $p-x$ preserves the direction:
\[
\frac{p-x}{x(1-x)} \;\geq\; 4(p-x) \qquad \text{for } x \in [q,p].
\]
Integrating over $[q,p]$,
\[
\int_q^p \frac{p-x}{x(1-x)}\, dx \;\geq\; 4 \int_q^p (p-x)\, dx \;=\; 2(p-q)^2.
\]

\textbf{Case $q > p$.} This is the \emph{same} computation with two sign reversals that cancel, so we do not repeat the algebra. First, reversing the limits, $\int_q^p = -\int_p^q$. Second, on $[p,q]$ we now have $x \geq p$, so $p-x \leq 0$, and multiplying $\frac{1}{x(1-x)} \geq 4$ by this nonpositive factor \emph{reverses} the pointwise inequality to $\frac{p-x}{x(1-x)} \leq 4(p-x)$. Integrating over $[p,q]$ and negating both sides (the second flip) restores the original direction:
\[
\int_p^q \frac{p-x}{x(1-x)}\, dx \;\leq\; 4\int_p^q (p-x)\, dx
\;\;\Longrightarrow\;\;
\int_q^p \frac{p-x}{x(1-x)}\, dx \;\geq\; 4\int_q^p (p-x)\, dx \;=\; 2(p-q)^2.
\]
Either way, $f_p(p) - f_p(q) \geq 2(p-q)^2$. \qedhere

\medskip
\emph{Remark (one-shot version).} Equivalently, note $f_p'(x) - 4(p-x) = (p-x)\big[\tfrac{1}{x(1-x)} - 4\big]$ has the same sign as $p-x$, since the bracket is $\geq 0$ on $(0,1)$. Hence the integrand is $\geq 0$ on $[q,p]$ when $q \leq p$ and $\leq 0$ on $[p,q]$ when $q > p$, so $\int_q^p \big[f_p'(x) - 4(p-x)\big]\, dx \geq 0$ in both cases.

\medskip
\emph{Alternative (no case split).} The same FTC computation avoids the case analysis entirely via a linear change of variables along the segment from $q$ to $p$. Using $2(p-q)^2 = 4\int_q^p (p-x)\, dx$,
\[
f_p(p) - f_p(q) - 2(p-q)^2 = \int_q^p \big[\,f_p'(x) - 4(p-x)\,\big]\, dx.
\]
Substitute $x = q + t(p-q)$ for $t \in [0,1]$, so $dx = (p-q)\, dt$ and $p - x = (1-t)(p-q)$:
\[
\int_q^p \big[\,f_p'(x) - 4(p-x)\,\big]\, dx
= \int_0^1 (1-t)\,(p-q)^2 \Big[\tfrac{1}{x(1-x)} - 4\Big]\, dt
\;\geq\; 0,
\]
since $(1-t) \geq 0$, $(p-q)^2 \geq 0$, and $\tfrac{1}{x(1-x)} - 4 \geq 0$ because $x = q + t(p-q) \in [0,1]$ is a convex combination of $q, p \in [0,1]$, so $x(1-x) \leq \tfrac14$. The orientation issue disappears because the $(p-x)$ in the integrand and the $(p-q)$ from $dx$ combine into $(p-q)^2 \geq 0$.
\end{solution}

%===================================
\section{Cumulative prediction error bound}\label{prob:bound}

Let $\mu$ be any environment, and assume it generates our actual universe.

\begin{definition}[Cumulative expected squared prediction error]\label{def:S_infty}
For environments / predictors $\nu, \rho$ and horizon $n \geq 1$, define the \textbf{cumulative expected squared error} as
\begin{align*}
\Sn{\nu}{\rho} ~&:=~ \sum_{t=1}^{n} \mathbb{E}_{X_{<t} \sim \nu}\!\Bigl[\,\sum_{x_t \in \mathbb{B}} \bigl(\nu(x_t \mid X_{<t}) - \rho(x_t \mid X_{<t})\bigr)^2\Bigr] \\
                &\phantom{:}=~ \sum_{t=1}^n \sum_{x_{<t} \in \mathbb{B}^{t-1}} \nu(x_{<t}) \sum_{x_t \in \mathbb{B}} \bigl(\nu(x_t \mid x_{<t}) - \rho(x_t \mid x_{<t})\bigr)^2.
\end{align*}
The first argument always names the distribution the expectation is taken over. Write $S^{\mu}_{\infty}(P) := \lim_{n \to \infty} \Sn{\mu}{P}$ for the infinite-horizon version with the true environment $\mu$ in the first slot; when $P = \xi$ (the Bayesian mixture of \cref{def:mixture}) we simply write $S^{\mu}_{\infty}$.
\end{definition}

\mysubsection{1}{[20] (Cumulative prediction bound)\label{prob:bound_ex}}

Show that
\[
S_{\infty}^{\mu} \leq - \ln w_\mu.
\]

\textit{Hint:} Apply \cref{thm:pinsker} to each summand to pass from squared error to KL divergence, then make the infinite sum explicit as a limit of finite sums up to $n$, then use \cref{prob:telescope} to telescope the KL divergences, and finally use the mixture representation of $\xi$.

\begin{solution}
The proof is a four-step chain:
\[
  S^\mu_\infty
  ~\stackrel{(1)}{\leq}~
  \sum_{t=1}^\infty \dt{\mu}{\xi}
  ~\stackrel{(2)}{=}~
  \lim_{n \to \infty} \Dn{\mu}{\xi}
  ~\stackrel{(3)}{\leq}~
  \lim_{n \to \infty} \sum_{x_{1:n}} \mu(x_{1:n}) \,\ln \tfrac{1}{w_\mu}
  ~\stackrel{(4)}{=}~
  -\ln w_\mu.
\]

\medskip
\textit{Step (1): Pinsker, pointwise.}
For each $t$ and $x_{<t}$, both $\mu(\cdot \mid x_{<t})$ and $\xi(\cdot \mid x_{<t})$ are probability distributions on $\mathbb{B}$ ($\xi$ by \cref{prob:one_step}), so \cref{thm:pinsker} gives
\[
\sum_{x_t} \bigl(\xi(x_t \mid x_{<t}) - \mu(x_t \mid x_{<t})\bigr)^2 \leq \sum_{x_t} \mu(x_t \mid x_{<t}) \,\ln \tfrac{\mu(x_t \mid x_{<t})}{\xi(x_t \mid x_{<t})}.
\]
Multiplying by $\mu(x_{<t}) \geq 0$ and summing over $x_{<t}$ and $t$ converts the LHS into $S^\mu_\infty$ and the RHS into $\sum_t \dt{\mu}{\xi}$ by \cref{def:Dn}.

\medskip
\textit{Step (2): Telescoping.}
By \cref{prob:telescope}, $\Dn{\mu}{\xi} = \sum_{t=1}^n \dt{\mu}{\xi}$. Taking $n \to \infty$ gives the equality.

\medskip
\textit{Step (3): Mixture domination.}
By the mixture representation \eqref{eq:mixture_expression},
\[
\xi(x_{1:n}) = \sum_{\nu \in \Mc} w_\nu\, \nu(x_{1:n}) \geq w_\mu\, \mu(x_{1:n}),
\]
so $\mu(x_{1:n}) / \xi(x_{1:n}) \leq 1/w_\mu$, hence $\ln(\mu/\xi) \leq \ln(1/w_\mu)$ pointwise. Substituting into the definition $\Dn{\mu}{\xi} = \sum_{x_{1:n}} \mu(x_{1:n}) \ln(\mu/\xi)$ gives the bound.

\medskip
\textit{Step (4): Total mass.}
Pull the constant out and use $\sum_{x_{1:n}} \mu(x_{1:n}) = 1$ ($\mu$ is an environment). \qedhere
\end{solution}

%===================================
\section{Explicit complexity bound}\label{prob:bound_explicit}

\mysubsection{1}{[05] (Explicit complexity bound)\label{prob:bound_explicit_ex}}

Specialize \cref{prob:bound} to the \textbf{Solomonoff prior} $w_\nu = 2^{-K(\nu)}$ (where $K(\nu)$ is the length of the shortest program computing $\nu$) to show
\[
S_{\infty}^{\mu} \leq K(\mu) \ln 2.
\]

\textit{Hint:} Substitute $w_\mu = 2^{-K(\mu)}$ into the bound from \cref{prob:bound}.

\begin{solution}
By \cref{prob:bound}, for \emph{any} prior, $S_{\infty}^{\mu} \leq -\ln w_\mu$. The Solomonoff prior assigns $\mu$ the weight $w_\mu = 2^{-K(\mu)}$, so
\[
S_{\infty}^{\mu} ~\leq~ -\ln w_\mu ~=~ -\ln 2^{-K(\mu)} ~=~ K(\mu) \ln 2. \qedhere
\]
\end{solution}


%===================================
\section*{Interpretation}

\cref{prob:bound,prob:bound_explicit} show that the \emph{cumulative} expected prediction error $S_\infty^\mu$ is finite. Since $S_\infty^\mu$ is a sum of non-negative terms, this immediately implies that the per-step expected prediction error
\[
d_t := \sum_{x_{<t} \in \mathbb{B}^*} \mu(x_{<t}) \sum_{x_t \in \mathbb{B}} \big( \xi(x_t \mid x_{<t}) - \mu(x_t \mid x_{<t}) \big)^2
\]
converges to zero as $t \to \infty$. In other words, $\xi$'s predictions become indistinguishable from $\mu$'s predictions on average over histories.

%===================================
\section{Universal dominance}\label{prob:dominance}

The mixture $\xi$ never assigns much less probability than any single environment
weighted by its prior. Under the Solomonoff prior this turns into a quantitative
``every computable environment is within a factor $2^{-K(\nu)}$ of $\xi$''
statement --- which is exactly the content of \cref{prob:bound,prob:bound_explicit}
viewed at a single time-step.

\mysubsection{1}{[05] (Universal dominance)\label{prob:dominance_ex}}

Show that for every $\nu \in \Mc$ and every $x \in \mathbb{B}^*$,
\[
  \xi(x) ~\geq~ w_\nu \cdot \nu(x).
\]

\begin{solution}
By \cref{def:mixture}, $\xi(x) = \sum_{\nu' \in \Mc} w_{\nu'} \,\nu'(x)$.
Every term in this sum is non-negative, so dropping all terms except the one with
$\nu' = \nu$ can only decrease the sum:
\[
  \xi(x) ~=~ \sum_{\nu' \in \Mc} w_{\nu'}\, \nu'(x)
  ~\geq~ w_\nu\, \nu(x). \qedhere
\]
\end{solution}

\begin{remark}[Solomonoff specialization]
Under the Solomonoff prior, $w_\nu = 2^{-K(\nu)}$, so the bound becomes
\[
  \xi_U(x) ~\geq~ 2^{-K(\nu)} \cdot \nu(x) \qquad \text{for every computable }\nu.
\]
This is what makes $\xi_U$ \emph{universal}: a single predictor dominates the
entire computable model class up to a factor that depends only on the model's
description length.
\end{remark}

\begin{remark}
This single inequality is the engine behind \cref{prob:bound}: dividing through
gives $\mu(x)/\xi(x) \leq 1/w_\mu$, which is then fed into Pinsker and the chain
rule in the proof of the main cumulative bound.
\end{remark}

%===================================
\section{From cumulative bound to a step-count}\label{prob:error_count}

\Cref{prob:bound,prob:bound_explicit} give an \emph{asymptotic} guarantee: the total
expected squared prediction error is finite. We can turn this directly into a
\emph{finite} statement: ``after at most this many steps, the mixture has
predicted to within $\varepsilon$ on average.''

\mysubsection{1}{[05] (Cumulative bound to step-count)\label{prob:error_count_ex}}

Define the per-step expected squared error
\[
  d_t ~:=~ \sum_{x_{<t} \in \mathbb{B}^*} \mu(x_{<t})
    \sum_{x_t \in \mathbb{B}} \big(\xi(x_t \mid x_{<t}) - \mu(x_t \mid x_{<t})\big)^2,
  \qquad \text{so } S^\mu_\infty = \sum_{t=1}^\infty d_t.
\]
Show that for every $\varepsilon > 0$,
\[
  \#\big\{\, t \geq 1 \;:\; d_t > \varepsilon \,\big\}
  ~<~ \frac{-\ln w_\mu}{\varepsilon}.
\]

\textit{Hint.} If $N$ such steps exist, what does that say about
$N\varepsilon$ compared to $\sum_t d_t$?

\begin{solution}
Let $N := \#\{t : d_t > \varepsilon\}$. Each bad step contributes more than
$\varepsilon$ to the cumulative sum, while all other steps contribute
non-negatively, so
\[
  N \cdot \varepsilon ~<~ \sum_{t \, : \, d_t > \varepsilon} d_t
  ~\leq~ \sum_{t=1}^\infty d_t ~=~ S^\mu_\infty ~\leq~ -\ln w_\mu,
\]
using \cref{prob:bound} at the last step. Dividing by $\varepsilon$ gives
$N < -\ln w_\mu / \varepsilon$. \qedhere
\end{solution}

\begin{remark}[Solomonoff specialization]
Under the Solomonoff prior $w_\mu = 2^{-K(\mu)}$, so $-\ln w_\mu = K(\mu)\ln 2$
and the bound becomes
\[
  \#\big\{\, t \geq 1 \;:\; d_t > \varepsilon \,\big\} ~<~ \frac{K(\mu)\ln 2}{\varepsilon}.
\]
\end{remark}

\begin{remark}
Strictly fewer than $-\ln w_\mu / \varepsilon$ steps can ever see per-step
expected error above $\varepsilon$; every other step is at most $\varepsilon$-bad
on average. The bound is independent of the time-step $t$ and of any structure of
the sequence beyond the prior weight of the true environment.
\end{remark}

%===================================
\section{Bounded number of prediction mistakes}\label{prob:mistakes}

So far we have bounded \emph{squared error}. For a deterministic environment
--- i.e.\ a single infinite binary sequence $x_1^*, x_2^*, \dots$ in $\Mc$ ---
one can ask the sharper $0/1$-loss question: how often does the mixture predict
the \emph{wrong bit}? The cumulative bound is strong enough to give an answer.

Let $\mu \in \Mc$ be a deterministic environment, so $\mu(x_t \mid x_{<t}) \in
\{0,1\}$ on every history; write $x_t^*$ for the unique bit on which
$\mu(\cdot \mid x_{<t})$ puts all its mass. Define the \textbf{threshold
predictor} associated with $\xi$: at each step, predict the more-likely bit,
\[
  \hat{x}_t ~:=~ \arg\max_{b \in \mathbb{B}} \xi(b \mid x_{<t})
  \qquad \text{(break ties arbitrarily).}
\]
We say the mixture makes a \textbf{mistake} at time $t$ if $\hat{x}_t \neq x_t^*$.

\mysubsection{1}{[15] (Bounded prediction mistakes)\label{prob:mistakes_ex}}

Show that the number of mistakes made by the threshold predictor on the
$\mu$-trajectory is at most
\[
  \#\{\, t \geq 1 \;:\; \hat{x}_t \neq x_t^* \,\}
  ~\leq~ -2 \ln w_\mu.
\]

\textit{Hint.} At a mistake step, $\xi(x_t^* \mid x_{<t}) \leq \tfrac12$ (why?).
The per-step squared error along the deterministic $\mu$-trajectory
simplifies dramatically; bound it from below, then sum.

\begin{solution}
\textbf{Step 1: per-step squared error along the $\mu$-trajectory.}
Because $\mu$ is deterministic, the outer sum over $x_{<t}$ in $d_t$ collapses
onto the unique trajectory $x_{<t}^* := x_1^* \cdots x_{t-1}^*$, with weight
$\mu(x_{<t}^*) = 1$. Write $p_t := \xi(x_t^* \mid x_{<t}^*) \in [0,1]$ for the
mass the mixture places on the \emph{correct} bit. Since $\mu$ assigns mass $1$
to $x_t^*$ and $0$ to the other bit,
\[
  d_t ~=~ \sum_{x_t \in \mathbb{B}} \big(\xi(x_t \mid x_{<t}^*) - \mu(x_t \mid x_{<t}^*)\big)^2
  ~=~ (p_t - 1)^2 + (1 - p_t - 0)^2 ~=~ 2(1 - p_t)^2.
\]

\textbf{Step 2: a mistake step contributes at least $\tfrac12$.}
A mistake at time $t$ means $\hat{x}_t \neq x_t^*$, i.e.\ the threshold predictor
chose the \emph{other} bit. This requires
$\xi(\hat{x}_t \mid x_{<t}^*) \geq \xi(x_t^* \mid x_{<t}^*)$, i.e.\
$1 - p_t \geq p_t$, i.e.\ $p_t \leq \tfrac12$. So at every mistake step,
\[
  d_t ~=~ 2(1 - p_t)^2 ~\geq~ 2 \cdot \tfrac{1}{4} ~=~ \tfrac12.
\]

\textbf{Step 3: sum and apply the cumulative bound.}
Let $N$ be the total number of mistakes. Summing $d_t$ over mistake steps
(non-negative contributions from non-mistake steps only help):
\[
  \tfrac12 \cdot N
  ~\leq~ \sum_{t \,:\, \hat{x}_t \neq x_t^*} d_t
  ~\leq~ \sum_{t=1}^\infty d_t ~=~ S^\mu_\infty
  ~\leq~ -\ln w_\mu,
\]
where the last step is \cref{prob:bound}. Rearranging, $N \leq -2 \ln w_\mu$.
\qedhere
\end{solution}

\begin{remark}[Solomonoff specialization]
Under the Solomonoff prior $w_\mu = 2^{-K(\mu)}$, so $-\ln w_\mu = K(\mu)\ln 2$
and the bound becomes
\[
  \#\{\, t \geq 1 \;:\; \hat{x}_t \neq x_t^* \,\} ~\leq~ 2 K(\mu) \ln 2.
\]
With $\ln 2 \approx 0.693$ the constant is just under $1.4$, so an environment
that takes $k$ bits to describe costs at most about $1.4 k$ mistakes \emph{ever}.
\end{remark}

\begin{remark}
This is the headline consequence of the cumulative bound: on \emph{any}
in-class deterministic infinite binary sequence, the mixture's threshold
predictor makes only finitely many mistakes, with the total proportional to
$-\ln w_\mu$. The bound is purely existence-style: it does not say \emph{when}
the mistakes happen. They could all occur at the start, or be spread out
unpredictably --- the bound only guarantees a finite total.
\end{remark}

%===================================
\section{Pareto optimality under KL loss}\label{prob:kl}

The final two sections establish that the Bayesian mixture $\xi$ is
\emph{uniquely Pareto-optimal} under both KL and squared prediction loss ---
no other predictor can do at least as well as $\xi$ on every environment
$\nu \in \Mc$ without coinciding with $\xi$ entirely. Hutter~\cite[\S3.6]{Hutter:04uaibook} states a weaker (non-uniqueness) version of this result for several losses including the two we treat here; the Pythagorean-identity proofs we use give the stronger uniqueness statement essentially for free.

A \textbf{predictor} (in this part) is any function $\rho$ assigning to each
history $x_{<t}$ a distribution $\rho(\cdot \mid x_{<t}) \in \Delta\mathbb{B}$;
by chain rule, $\rho(x_{1:n}) := \prod_t \rho(x_t \mid x_{<t})$. The mixture
$\xi$ is one such predictor, and so is any $\nu \in \Mc$.

\begin{definition}[Pareto domination]\label{def:pareto}
For a loss $\mathcal F(\nu, \rho)$ that takes an environment $\nu$ and a
predictor $\rho$, a predictor $\rho$ \textbf{weakly Pareto-dominates} $\xi$ if
\[
  \mathcal F(\nu, \rho) ~\leq~ \mathcal F(\nu, \xi)
  \qquad \text{for every } \nu \in \Mc.
\]
$\xi$ is \textbf{Pareto-optimal} (w.r.t.\ $\mathcal F$) if the only predictor
that weakly dominates it is $\xi$ itself.\footnote{This is stronger than the
usual ``no strict domination'' formulation: weak domination already forces
$\rho = \xi$, so a fortiori there is no $\rho$ that is strictly better on some
$\nu$ without being worse on another.} The quantifier ranges over the whole
class --- there is no single ``true environment.'' Each $\nu$ takes its turn
as the candidate environment; Pareto-optimality says no $\rho$ beats $\xi$
uniformly across the class.

We specialize $\mathcal F$ to $\Dn{\nu}{\rho}$ in \cref{prob:kl} below (KL loss, joint level) and to $\Sn{\nu}{\rho}$ in \cref{prob:sq} (cumulative squared loss).
\end{definition}

KL is the easier of the two cases. It lives entirely at the joint distribution
level $\nu(x_{1:n})$, where the Bayesian mixing identity $\sum_\nu w_\nu \,
\nu(x_{1:n}) = \xi(x_{1:n})$ is literally the definition of $\xi$. The
Pythagorean decomposition and the Pareto aggregation use exactly the same
weights, and the proof is a one-step calculation.

\mysubsection{1}{[10] (KL Pythagorean identity)\label{prob:kl_pythag}}

For \emph{any} predictor $\rho$ with $\rho(x_{1:n}) > 0$ on every $x_{1:n} \in \mathbb{B}^n$, show that
\[
  \Dn{\xi}{\rho} ~+~ \sum_{\nu \in \Mc} w_\nu \, \Dn{\nu}{\xi}
  ~=~ \sum_{\nu \in \Mc} w_\nu \, \Dn{\nu}{\rho}.
\]
This is an equality of three real numbers; there is no history variable
floating free --- all references to histories $x_{1:n}$ are summed over inside
the $D_n$'s, and weighted by the relevant joint distribution (per \cref{def:Dn}).

\textit{Hint.} Subtract the second sum from the third, simplify the
log-difference $\ln\tfrac{\nu}{\rho} - \ln\tfrac{\nu}{\xi} = \ln\tfrac{\xi}{\rho}$
inside the $\nu$-expectation, then swap the order of summation and use the
prior-mixing identity $\sum_\nu w_\nu \nu(x_{1:n}) = \xi(x_{1:n})$ (which holds
by definition of $\xi$ as the prior mixture on joints).

\begin{solution}
Take the difference $\sum_\nu w_\nu \Dn{\nu}{\rho} - \sum_\nu w_\nu \Dn{\nu}{\xi}$.
Inside each $D_n$ the $\nu(x_{1:n}) \ln \nu(x_{1:n})$ pieces cancel, leaving
\begin{align*}
\sum_\nu w_\nu \Dn{\nu}{\rho} - \sum_\nu w_\nu \Dn{\nu}{\xi}
&= \sum_\nu w_\nu \sum_{x_{1:n}}
   \nu(x_{1:n})\,\Bigl[\ln\tfrac{\nu(x_{1:n})}{\rho(x_{1:n})}
                  - \ln\tfrac{\nu(x_{1:n})}{\xi(x_{1:n})}\Bigr] \\
&= \sum_\nu w_\nu \sum_{x_{1:n}}
   \nu(x_{1:n})\,\ln \tfrac{\xi(x_{1:n})}{\rho(x_{1:n})}.
\end{align*}
The factor $\ln(\xi/\rho)$ does not depend on $\nu$, so we swap the order of
summation and pull it out, then use $\sum_\nu w_\nu \nu(x_{1:n}) = \xi(x_{1:n})$:
\begin{align*}
\sum_\nu w_\nu \sum_{x_{1:n}}
   \nu(x_{1:n})\,\ln \tfrac{\xi(x_{1:n})}{\rho(x_{1:n})}
&= \sum_{x_{1:n}} \ln \tfrac{\xi(x_{1:n})}{\rho(x_{1:n})}
   \underbrace{\sum_\nu w_\nu \nu(x_{1:n})}_{=\, \xi(x_{1:n})} \\
&= \sum_{x_{1:n}} \xi(x_{1:n}) \,\ln \tfrac{\xi(x_{1:n})}{\rho(x_{1:n})}
~=~ \Dn{\xi}{\rho}.
\end{align*}
Rearranging gives the claimed identity. \qedhere
\end{solution}

\mysubsection{2}{[05] (KL Pareto-optimality)\label{prob:kl_pareto}}

Show that $\xi$ is Pareto-optimal under $\Dn{\nu}{\rho}$ in the sense of \cref{def:pareto}: if $\rho$ is any predictor with
\[
  \Dn{\nu}{\rho} ~\leq~ \Dn{\nu}{\xi} \qquad \text{for every } \nu \in \Mc,
\]
then $\rho(x_{1:n}) = \xi(x_{1:n})$ for \emph{every} $x_{1:n} \in \mathbb{B}^n$
(and hence, by chain rule, $\rho(x_t \mid x_{<t}) = \xi(x_t \mid x_{<t})$ at
every history $x_{<t}$ with $\xi(x_{<t}) > 0$).

\textit{Hint.} Multiply the assumed inequality by $w_\nu \geq 0$ and sum over
$\nu \in \Mc$; then apply the KL Pythagorean identity to recognize the extra
non-negative term --- which must vanish.

\begin{solution}
Multiply the assumed inequality by $w_\nu \geq 0$ and sum over $\nu \in \Mc$:
\begin{equation}\label{eq:kl_weighted_dom}
  \sum_\nu w_\nu\, \Dn{\nu}{\rho} ~\leq~ \sum_\nu w_\nu\, \Dn{\nu}{\xi}.
\end{equation}
The KL Pythagorean identity rewrites the LHS as
\[
  \Dn{\xi}{\rho} \;+\; \sum_\nu w_\nu\, \Dn{\nu}{\xi},
\]
so \cref{eq:kl_weighted_dom} forces $\Dn{\xi}{\rho} \leq 0$. KL is non-negative, hence $\Dn{\xi}{\rho} = 0$, which forces $\rho(x_{1:n}) = \xi(x_{1:n})$ for every $x_{1:n} \in \mathbb{B}^n$. By chain rule, the conditionals agree at every history $x_{<t}$ with $\xi(x_{<t}) > 0$. \qedhere
\end{solution}

%===================================
\section{Pareto optimality under squared loss}\label{prob:sq}

The squared specialization $\mathcal F(\nu, \rho) = \Sn{\nu}{\rho}$ (\cref{def:S_infty}) lives one level down from KL: at the conditional distributions
$\nu(\cdot \mid x_{<t})$. The natural Pythagorean decomposition at a fixed
history uses \emph{posterior} weights $w(\nu \mid x_{<t})$ --- those are the
weights that make $\xi(x_t \mid x_{<t})$ the mean of the $\nu(x_t \mid x_{<t})$'s
(via the posterior-predictive form, \cref{def:mixture}). The Pareto-optimality
aggregation, however, uses prior weights $w_\nu$. Bridging the two takes one
extra Bayes-rule step.

\mysubsection{1}{[10] (Squared Pythagorean identity)\label{prob:sq_pythag}}

Fix any $t \in \{1, \dots, n\}$ and any history $x_{<t} \in \mathbb{B}^{t-1}$
--- \emph{the history is arbitrary, not sampled from any environment}. For any
$x_t \in \mathbb{B}$ and any predictor $\rho$, show
\begin{align*}
  &\bigl(\xi(x_t \mid x_{<t}) - \rho(x_t \mid x_{<t})\bigr)^{2}
  ~+~
  \sum_{\nu \in \Mc} w(\nu \mid x_{<t})\,\bigl(\nu(x_t \mid x_{<t}) - \xi(x_t \mid x_{<t})\bigr)^{2} \\
  &=~
  \sum_{\nu \in \Mc} w(\nu \mid x_{<t})\,\bigl(\nu(x_t \mid x_{<t}) - \rho(x_t \mid x_{<t})\bigr)^{2}.
\end{align*}

\textit{Hint.} Add and subtract $\xi(x_t \mid x_{<t})$ inside the squared term,
then expand. Use $\sum_\nu w(\nu \mid x_{<t}) = 1$ and the posterior-predictive
form of $\xi$ (\cref{def:mixture}), $\sum_\nu w(\nu \mid x_{<t})\,\nu(x_t \mid x_{<t})
= \xi(x_t \mid x_{<t})$, to kill the cross-term. The weighting must be
\emph{posterior} weights $w(\nu \mid x_{<t})$ (not prior weights $w_\nu$),
because $\xi(x_t \mid x_{<t})$ is the posterior-weighted mean of the
$\nu(x_t \mid x_{<t})$'s --- not the prior one.

\begin{solution}
Throughout, the history $x_{<t}$ and symbol $x_t$ are fixed but arbitrary.
Abbreviate $\nu := \nu(x_t \mid x_{<t})$, $\xi := \xi(x_t \mid x_{<t})$,
$\rho := \rho(x_t \mid x_{<t})$, and $\tilde w_\nu := w(\nu \mid x_{<t})$ for
the duration of the calculation.

\medskip
\textit{Step 1: $\xi$ is the posterior-weighted mean of the $\nu$'s.}
By the posterior-predictive form of the mixture (\cref{def:mixture}) and the
posterior normalization $\sum_\nu \tilde w_\nu = 1$,
\[
  \sum_\nu \tilde w_\nu \, \nu ~=~ \xi.
\]

\medskip
\textit{Step 2: the cross-term vanishes.}
A direct consequence of Step 1, again using $\sum_\nu \tilde w_\nu = 1$:
\[
  \sum_\nu \tilde w_\nu \, (\nu - \xi)
  ~=~ \sum_\nu \tilde w_\nu \, \nu \;-\; \xi \sum_\nu \tilde w_\nu
  ~=~ \xi - \xi ~=~ 0.
\]

\medskip
\textit{Step 3: expand the square.}
Write $\nu - \rho = (\nu - \xi) + (\xi - \rho)$ and expand:
\begin{align*}
\sum_\nu \tilde w_\nu (\nu - \rho)^2
  &= \sum_\nu \tilde w_\nu\bigl[(\nu - \xi)^2 + 2(\nu - \xi)(\xi - \rho) + (\xi - \rho)^2\bigr] \\
  &= \sum_\nu \tilde w_\nu (\nu - \xi)^2
     \;+\; 2(\xi - \rho)\sum_\nu \tilde w_\nu (\nu - \xi)
     \;+\; (\xi - \rho)^2 \sum_\nu \tilde w_\nu.
\end{align*}
The middle sum is zero by Step 2; the trailing $\sum_\nu \tilde w_\nu = 1$. So
\[
  (\xi - \rho)^2 \;+\; \sum_\nu \tilde w_\nu (\nu - \xi)^2
  ~=~ \sum_\nu \tilde w_\nu (\nu - \rho)^2. \qedhere
\]
\end{solution}

\mysubsection{2}{[05] (Squared Pareto-optimality)\label{prob:sq_pareto}}

Show that $\xi$ is Pareto-optimal under $\Sn{\nu}{\rho}$ in the sense of \cref{def:pareto}: if $\rho$ is any predictor with
\[
  \Sn{\nu}{\rho} ~\leq~ \Sn{\nu}{\xi} \qquad \text{for every } \nu \in \Mc,
\]
then $\rho(x_t \mid x_{<t}) = \xi(x_t \mid x_{<t})$ at every history $x_{<t} \in
\mathbb{B}^{t-1}$ \emph{with $\xi(x_{<t}) > 0$} (equivalently, $\xi$-almost
surely) and every $x_t \in \mathbb{B}$. Histories that the mixture never reaches
($\xi(x_{<t}) = 0$, i.e., reached by no $\nu \in \Mc$ either) are invisible
to the loss and are not pinned down.

\textit{Hint.} Multiply the assumed inequality by $w_\nu \geq 0$ and sum over
$\nu$; then for each fixed history $x_{<t}$, multiply the per-history
Pythagorean identity from above by $\xi(x_{<t})$ and use the Bayes identity
$w(\nu \mid x_{<t})\,\xi(x_{<t}) = w_\nu \,\nu(x_{<t})$ to bridge from posterior
weights (inside the identity) to prior weights (in the aggregated Pareto
inequality).

\begin{solution}
Multiply the assumed inequality by $w_\nu \geq 0$ and sum over $\nu \in \Mc$:
\begin{equation}\label{eq:sq_weighted_dom}
  \sum_\nu w_\nu\,\Sn{\nu}{\rho}
  ~\leq~
  \sum_\nu w_\nu\,\Sn{\nu}{\xi}.
\end{equation}

Now bridge to the per-history Pythagorean. At each \emph{fixed} history $x_{<t}
\in \mathbb{B}^{t-1}$ and symbol $x_t \in \mathbb{B}$, the identity says
\[
  (\xi - \rho)^2 \;+\; \sum_\nu w(\nu \mid x_{<t})\,(\nu - \xi)^2
  ~=~ \sum_\nu w(\nu \mid x_{<t})\,(\nu - \rho)^2,
\]
with the abbreviations $\nu := \nu(x_t \mid x_{<t})$, $\xi := \xi(x_t \mid x_{<t})$,
$\rho := \rho(x_t \mid x_{<t})$. Multiply by $\xi(x_{<t})$ and use Bayes,
$w(\nu \mid x_{<t})\,\xi(x_{<t}) = w_\nu \,\nu(x_{<t})$:
\[
  \xi(x_{<t})(\xi - \rho)^2
  \;+\; \sum_\nu w_\nu \,\nu(x_{<t})\,(\nu - \xi)^2
  ~=~ \sum_\nu w_\nu \,\nu(x_{<t})\,(\nu - \rho)^2.
\]
Sum over $t = 1, \dots, n$ and all histories $x_{<t} \in \mathbb{B}^{t-1}$ and
symbols $x_t \in \mathbb{B}$. The two $\sum_\nu w_\nu \,\nu(x_{<t})(\cdots)$
terms become exactly $\sum_\nu w_\nu \,\Sn{\nu}{\xi}$ and $\sum_\nu w_\nu \,
\Sn{\nu}{\rho}$ respectively, so
\[
  \|\xi - \rho\|_\xi^2
  \;+\; \sum_\nu w_\nu\,\Sn{\nu}{\xi}
  ~=~
  \sum_\nu w_\nu\,\Sn{\nu}{\rho},
\]
where
\[
  \|\xi - \rho\|_\xi^2 ~:=~
  \sum_{t=1}^n \sum_{x_{<t}} \xi(x_{<t}) \sum_{x_t}
    \bigl(\xi(x_t \mid x_{<t}) - \rho(x_t \mid x_{<t})\bigr)^2 ~\geq~ 0.
\]
Combining with \cref{eq:sq_weighted_dom},
\[
  \|\xi - \rho\|_\xi^2 \;+\; \sum_\nu w_\nu\,\Sn{\nu}{\xi}
  ~\leq~ \sum_\nu w_\nu\,\Sn{\nu}{\xi}
  \;\;\Longrightarrow\;\;
  \|\xi - \rho\|_\xi^2 ~\leq~ 0.
\]
Since $\|\xi - \rho\|_\xi^2$ is a sum of non-negative terms and is itself
$\leq 0$, every term vanishes:
$\xi(x_{<t})\bigl(\xi(x_t \mid x_{<t}) - \rho(x_t \mid x_{<t})\bigr)^2 = 0$
at every $(t, x_{<t}, x_t)$. So $\rho(x_t \mid x_{<t}) = \xi(x_t \mid x_{<t})$
at every history $x_{<t}$ with $\xi(x_{<t}) > 0$. \qedhere
\end{solution}

%===================================
\section*{Closing remarks}

\begin{remark}[Why these are ``Pythagorean'']
Both identities have the shape
\[
  \text{(loss from $\rho$, averaged over $\nu$'s)}
  \;=\;
  \text{(loss from $\xi$, averaged over $\nu$'s)} \;+\;
  \text{(loss from $\rho$ to $\xi$).}
\]
For squared loss this is exactly Pythagoras' theorem in Euclidean space at
each fixed history $x_{<t}$, with $\xi(\cdot \mid x_{<t})$ playing the role
of the orthogonal projection of the ``cloud of $\nu(\cdot \mid x_{<t})$'s''
onto the line of possible reference predictors. For KL it is the Bregman-divergence
analogue at the joint level $x_{1:n}$: $\KL$ behaves like a squared distance,
with $\xi(X_{1:n})$ as the Bregman projection. The crucial property both
proofs share is the ``mean-zero cross-term'': $\sum_\nu w(\nu \mid x_{<t})
(\nu - \xi) = 0$ at the conditional level (squared loss) and $\sum_\nu w_\nu
(\nu(x_{1:n}) - \xi(x_{1:n})) = 0$ at the joint level (KL). In each case the
weighting is whichever one makes $\xi$ the mean of the $\nu$'s --- posterior
weights for the conditional, prior weights for the joint.
\end{remark}

\begin{remark}[Why KL is easier than squared loss]
KL's Pythagorean identity uses prior weights $w_\nu$, matching the weights in
the Pareto sum directly: one identity, one inequality, done. The squared-loss
Pythagorean uses posterior weights $w(\nu \mid x_{<t})$, so bridging to the
prior-weighted Pareto sum costs an extra step (multiply by $\xi(x_{<t})$, apply
Bayes). Geometrically, KL lives at the joint $x_{1:n}$ level where one
identity does everything; squared loss lives at the conditional $(x_t \mid
x_{<t})$ level where the natural mean-zero weighting shifts with the history.
The conclusion is also slightly weaker for squared loss --- $\rho = \xi$ only
at $\xi$-positive histories (i.e., $\xi$-almost surely), not at every history
--- because off-support histories are invisible to a $\nu$-expected loss.
\end{remark}

\begin{remark}[$\xi$ minimizes the prior-weighted loss]
Both identities show more than Pareto optimality. For KL, reading the
joint-level identity as a function of $\rho$, the $\rho$-dependence on the
RHS is a single non-negative term $\Dn{\xi}{\rho}$ that vanishes only at
$\rho(x_{1:n}) = \xi(x_{1:n})$. So the unique minimizer of $\rho \mapsto
\sum_\nu w_\nu \,\Dn{\nu}{\rho}$ is $\rho = \xi$. For squared loss, after the
multiply-by-$\xi(x_{<t})$-and-sum aggregation, the $\rho$-dependence is
$\sum_{t,x_{<t},x_t} \xi(x_{<t})(\xi - \rho)^2$, which vanishes iff $\rho = \xi$
at every $\xi$-positive history. In both cases the Pareto-optimality result is
an immediate corollary.
\end{remark}

\begin{remark}[Versus the standard ``Bayes is admissible'' result]
Pareto optimality is a kind of admissibility statement: nothing weakly
dominates $\xi$ across the class. The standard textbook way to prove this for
Bayesian estimators goes via the complete-class theorem in decision theory and
is considerably more abstract. The Pythagorean identities used here give the
result by direct calculation, with no decision-theoretic machinery and no
assumptions beyond a proper prior ($\sum_\nu w_\nu = 1$). The reason it works
so cleanly is that both losses are convex Bregman divergences in their second
argument --- exactly the setting in which the Bayesian mixture is the unique
minimizer.
\end{remark}

%===================================
\section{Misspecified models: when $\mu \notin \Mc$}\label{prob:misspec}

The cumulative bound \cref{prob:bound} assumed the true environment $\mu$ lives
in the model class $\Mc$. What happens when it does not? We now show the
bound \emph{degrades gracefully}: it splits cleanly into a complexity term
(``cost of not knowing which model is best''), exactly as before, plus a
linear-in-time approximation term (``cost of no model being right''). See \cite[\S3.2.8]{Hutter:04uaibook} for the original treatment.

Throughout this section, $\mu$ is an arbitrary environment, possibly outside
$\Mc$. We will state the bound in terms of an arbitrary
$\hat\mu \in \Mc$ --- think of $\hat\mu$ as ``any in-class model you would
like to compare $\xi$ against.'' The bound holds simultaneously for every
choice of $\hat\mu \in \Mc$, so taking the infimum over $\hat\mu$ on the
right-hand side gives the tightest version.

\mysubsection{1}{[10] (Misspecified KL bound)\label{prob:misspec_ex}}

Show that for every $n \geq 1$ and every $\hat\mu \in \Mc$,
\[
  \Dn{\mu}{\xi} ~\leq~ -\ln w_{\hat\mu} ~+~ \Dn{\mu}{\hat\mu}.
\]

\textit{Hint.} Mimic the proof of \cref{prob:bound}, but bound $\xi$ from below
by $w_{\hat\mu} \,\hat\mu$ rather than $w_\mu \, \mu$.

\begin{solution}
Mixture dominance applied to $\hat\mu \in \Mc$ (\cref{prob:dominance}) gives
$\xi(x_{1:n}) \geq w_{\hat\mu}\, \hat\mu(x_{1:n})$ for every $x_{1:n}$. Hence
$\mu(x_{1:n}) / \xi(x_{1:n}) \leq \mu(x_{1:n}) / (w_{\hat\mu}\, \hat\mu(x_{1:n}))$, and taking logs,
\[
  \ln \frac{\mu(x_{1:n})}{\xi(x_{1:n})}
  ~\leq~ -\ln w_{\hat\mu} \;+\; \ln \frac{\mu(x_{1:n})}{\hat\mu(x_{1:n})}.
\]
Now take the $\mu$-expectation. The LHS becomes $\Dn{\mu}{\xi}$; the second RHS term becomes $\Dn{\mu}{\hat\mu}$; the constant $-\ln w_{\hat\mu}$ survives because $\mu(X_{1:n})$ has total mass $1$. \qedhere
\end{solution}

\begin{remark}[Reading the two terms]
The right-hand side splits into two qualitatively different terms:
\begin{itemize}\itemsep=2pt
\item $-\ln w_{\hat\mu}$ is \emph{constant in $n$} --- under the Solomonoff
prior, this is $K(\hat\mu)\ln 2$. The familiar ``complexity'' cost of search.
\item The approximation term $\Dn{\mu}{\hat\mu} = \sum_{t=1}^n \dt{\mu}{\hat\mu}$ (by \cref{prob:telescope}) is a sum of $n$ per-step KLs; in general it grows with $n$. It vanishes identically iff $\hat\mu$ matches $\mu$ on every $\mu$-reachable history --- in particular when $\mu \in \Mc$ and we pick $\hat\mu = \mu$, recovering \cref{prob:bound}.
\end{itemize}
Combining with Pinsker (\cref{thm:pinsker}) gives the corresponding bound on
$S^\mu_\infty$, which is infinite in general but inherits the rate of $\hat\mu$.
\end{remark}

\begin{remark}[What does ``best in class'' mean?]
The bound is valid for every $\hat\mu \in \Mc$. A tempting question: which
$\hat\mu$ gives the tightest bound? The answer is genuinely setting-dependent.

\textit{Finite horizon $n$.} The minimizer of the RHS is
\[
  \hat\mu_n ~:=~ \arg\min_{\nu \in \Mc}\Bigl\{ -\ln w_\nu \;+\; \Dn{\mu}{\nu} \Bigr\}.
\]
This is well-defined (finitely many terms, attained when $\Mc$ is finite),
but the answer depends on $n$ --- for small $n$ the prior term dominates and a
high-prior coarse model wins; for large $n$ the fit term dominates and the best
per-step approximator wins. Both regimes are correct.

\textit{Asymptotic rate.} A horizon-free analogue minimizes the per-step KL
rate $\lim_{n \to \infty} \tfrac{1}{n} \sum_{t \leq n} \dt{\mu}{\nu}$. This is the right notion when one cares about the linear-growth slope of the bound. But it has two drawbacks: the limit need not exist for non-stationary $\mu$ (one can substitute $\liminf$ but at the cost of clarity), and it discards the prior $w_\nu$ entirely, so the minimizer is not the same as $\hat\mu_n$ at any finite $n$ --- only of its slope.

\textit{Stationary $\mu$.} The per-step KLs $\dt{\mu}{\nu}$ are constant in $t$, so all three notions essentially agree: $\hat\mu = \arg\min_{\nu \in \Mc} D_{\text{KL}}(\mu \,\|\, \nu)$, where the KL is taken under the (common) one-step conditional law. This is the case where ``$\hat\mu$ = projection of $\mu$ onto $\Mc$ in KL'' has unambiguous meaning.

\textit{Upshot.} The bound holds for all $\hat\mu \in \Mc$, and one is free
to take the infimum on the RHS over $\hat\mu$. Which $\hat\mu$ actually
attains that infimum is a separate question whose answer depends on horizon,
prior, and any structural assumptions on $\mu$.
\end{remark}

%===================================
\clearpage
\appendix
\renewcommand{\thesection}{\Alph{section}}
\crefname{section}{Appendix}{Appendices}
\Crefname{section}{Appendix}{Appendices}

\section{Proof of Pinsker's inequality}\label{app:pinsker}

We reduce \cref{thm:pinsker} to the binary inequality \cref{lem:pinsker_binary}.

\begin{proof}[Proof of \cref{thm:pinsker}]
Recall $p_0 + p_1 = 1$ and $q_0 + q_1 = 1$. We handle the boundary case
$q_x = 0$ first, then reduce the interior case directly to
\cref{lem:pinsker_binary}.

\medskip
\textbf{Case 1: $q_x = 0$ for some $x \in \mathbb{B}$.}
If $p_x > 0$, then the right-hand side of $(\star)$ contains
$p_x \ln \tfrac{p_x}{0} = +\infty$, so $(\star)$ holds trivially.
If instead $p_x = 0$, that atom contributes $0$ to both sides
(using the convention $0 \ln \tfrac{0}{0} = 0$). The other atom $x'$ then
satisfies $p_{x'} = 1$ and $q_{x'} = 1 - q_x = 1$ --- so both sides vanish and
$(\star)$ holds.

\medskip
\textbf{Case 2: $q_0, q_1 > 0$.}
Since $q_1 = 1 - q_0$ and $p_1 = 1 - p_0$, the two sides of $(\star)$ become
\[
    \sum_{x \in \mathbb{B}} (q_x - p_x)^2
    = (q_0 - p_0)^2 + \big((1-q_0) - (1-p_0)\big)^2
    = 2(p_0 - q_0)^2,
\]
\[
    \sum_{x \in \mathbb{B}} p_x \ln \tfrac{p_x}{q_x}
    = p_0 \ln \tfrac{p_0}{q_0} + (1 - p_0) \ln \tfrac{1 - p_0}{1 - q_0}.
\]
So $(\star)$ is exactly
$2(p_0 - q_0)^2 \leq p_0 \ln \tfrac{p_0}{q_0} + (1 - p_0) \ln \tfrac{1 - p_0}{1 - q_0}$,
which is \cref{lem:pinsker_binary} with $p = p_0$ and $q = q_0$
(here $0 < q_0 < 1$, since $q_0, q_1 > 0$).
\end{proof}

%===================================
\bibliographystyle{alphaurl}
\bibliography{biblo}

\end{document}
